\documentclass[10pt,a4paper]{amsart}
\usepackage[margin=19mm]{geometry}
\usepackage{amsmath,amssymb,mathtools}
\usepackage[scr=rsfs,cal=euler]{mathalfa}
\usepackage{xfrac}
\usepackage[unicode,hidelinks,bookmarks=false]{hyperref}
\newcommand{\ie}{i.e.\ }
\newcommand{\cf}{cf.\ }
\newcommand{\resp}{respectively\ }
\newcommand{\Subsection}{\subsection}
\providecommand{\nobreakdash}{\nobreak\hbox{-}}
\setlength{\emergencystretch}{2em}
\multlinegap0pt
\usepackage{bm}
\usepackage{bbm}
\usepackage{dsfont}
\usepackage{xspace}
\usepackage{cancel}
\usepackage{mathtools}
\usepackage{amsthm}
\makeatletter
\@ifundefined{theorem}{\newtheorem{theorem}{Theorem}}{}
\makeatother
\newcommand{\G}{\mathbb{G}}
\newcommand{\M}[1]{\mathbb{#1}}
\newcommand{\mood}[1]{(\operatorname{mod} #1)}
\newcommand{\p}{\mathbb{P}}
\title[On The Infinitude of the Twin Primes]{On The Infinitude of the Twin Primes}
\author{Berndt Gensel}
\date{Source: arXiv:1909.07975v7 (CC BY 4.0)}
\begin{document}
\maketitle

\noindent\emph{Source and licence.} On The Infinitude of the Twin Primes. Version arXiv:1909.07975v7 (CC BY 4.0), distributed under the Creative Commons Attribution 4.0 International licence, as recorded in the arXiv licence metadata for this exact version and shown on the abstract page https://arxiv.org/abs/1909.07975v7. Full rights notice: licenses/arxiv-1909.07975-CC-BY-4.0.txt.\par\smallskip

\noindent\textit{Editorial scope.} Counted unit: the Theorem whose source label is \texttt{T-twinprimegenerators} in the article (source0.tex lines 136--143, source label \texttt{T-twinprimegenerators}), delivered together with the article's own complete original proof of it (source0.tex lines 144--195, one \texttt{proof} environment of 52 lines). No mathematics is written, repaired or re-derived: statement and proof are the article's own text line for line. Required context retained uncounted: none. Every definition, display and label of those lines is retained unchanged. Omitted from this package and not redistributed: 1--135, 196--427. Nothing inside a retained excerpt is omitted or altered: every retained body is delivered exactly as the article's lines are written, so no omission receipt of any kind accompanies this package. Citations: the retained excerpts carry no citation command at all, so no bibliography entry of the article is transcribed and no reference is redistributed. Reference adaptation: none was needed and none was made; every reference inside the delivered unit resolves inside it, so no unresolved reference or citation is produced. Printed slips kept literal: no printed word, symbol, hypothesis or number of the counted statement or of its proof was altered. Layout and class: the article's own class is \texttt{amsart} with packages including graphicx, geometry; the wrapper is the lane's pinned \texttt{amsart} editorial wrapper at 10pt on A4, which restates the theorem-like environments the retained text needs and adds the article's own macro definitions that the retained text needs (\texttt{\textbackslash G}, \texttt{\textbackslash M}, \texttt{\textbackslash mood}, \texttt{\textbackslash p}), with extra wrapper packages (bm, bbm, dsfont, xspace, cancel, mathtools). The delivered numbering is the unit's own. Assets: no retained span contains a figure environment, a graphics-inclusion command, a PDF-page inclusion, a table or a TikZ picture; no image file is copied into this package, the package declares no asset dependency and no image file is redistributed with it. Licence: version arXiv:1909.07975v7 of this article is distributed under the Creative Commons Attribution 4.0 International licence, as recorded in the arXiv licence metadata for this exact version and shown on the abstract page https://arxiv.org/abs/1909.07975v7. A full rights notice is delivered at licenses/arxiv-1909.07975-CC-BY-4.0.txt. Characters: every retained excerpt is pure ASCII, so the delivered engine, XeLaTeX with its default Latin Modern text font, reports no missing character.
\begin{theorem}\label{T-twinprimegenerators}
A number $x$ is a \textbf{twin prime generator}, a member of $\G$, if and only if there is \textbf{no} $p \in \p^*$ with $p<6x-1$ that one of the following congruences fulfills:
\begin{align}
x &\equiv - \kappa(p)\mood{p} \label{E-kappa-}
\\
x &\equiv + \kappa(p)\mood{p} \label{E-kappa+}
\end{align}
\end{theorem}

\begin{proof}
At first we assume, that there is a prime $p\in\p^*$ with $p<6x-1$ such that \eqref{E-kappa-} or \eqref{E-kappa+} is valid. There are two cases.
\begin{enumerate}
\item{
$p \in \p_-$ , which means $p=6\kappa(p)-1$:

If \eqref{E-kappa-} is true, then there is an $n \in \M{N}$ with
\begin{align*}
x &= -\kappa(p)+n \cdot (6\kappa(p)-1)\\
6x &= -6\kappa(p)+6n \cdot (6\kappa(p)-1)\\ 
6x+1 &=-6\kappa(p)+6n \cdot (6\kappa(p)-1)+1\\
&=(6n-1)(6\kappa(p)-1)\\
&=(6n-1)\cdot p\\
\implies 6x+1 &\equiv 0 \mood{p} \implies x \notin \G
\end{align*}
For \eqref{E-kappa+} the proof will be done with $6x-1$:
\begin{align*}
6x-1 &=6\kappa(p)+6n \cdot (6\kappa(p)-1)-1\\
&=(6n+1)(6\kappa(p)-1)\\
&=(6n+1)\cdot p\\
\implies 6x-1 &\equiv 0 \mood{p} \implies x \notin \G
\end{align*}
}
\item{
$p \in \p_+$, which means $p=6\kappa(p)+1$:

We proceed similarly with \eqref{E-kappa-} and $6x-1$ as well as \eqref{E-kappa+} and $6x+1$:
\begin{align*}
6x-1 &=(6n-1)(6\kappa(p)+1) \implies 6x-1 \equiv 0 \mood{p}\\
6x+1 &=(6n+1)(6\kappa(p)+1) \implies 6x+1 \equiv 0 \mood{p}
\end{align*}
}
\end{enumerate}
This implies $x \notin \G$, if the congruences \eqref{E-kappa-} or \eqref{E-kappa+} are valid. They cannot both be true, as they are mutually exclusive. \\

If on the other hand $x \notin \G$, then $6x-1$ or $6x+1$ is no prime. 
Let $6x-1 \equiv 0 \mood{p}$ for any $p \in \p_-$. Then we have
\begin{align*}
6x-1 &\equiv 0 \mood{p}\equiv p \mood{p}\\
&\equiv (6\kappa(p)-1) \mood{p}\\
6x &\equiv 6\kappa(p) \mood{p}\\
x &\equiv \kappa(p) \mood{p}\implies \eqref{E-kappa+}.
\end{align*}
For any $p \in \p_+$ we have
\begin{align*}
6x-1 &\equiv -p \mood{p}\\
&\equiv -(6\kappa(p)+1) \mood{p}\\
6x &\equiv -6\kappa(p) \mood{p}\\
x &\equiv -\kappa(p) \mood{p}\implies \eqref{E-kappa-}.
\end{align*}
The other both cases we can handle in the same way. Therefore either \eqref{E-kappa-} or \eqref{E-kappa+} is valid if $x \notin \G$.
\end{proof}

\end{document}
